% !TeX program = xelatex
\documentclass[12pt,a4paper]{article}

\usepackage[a4paper,margin=24mm]{geometry}
\usepackage{fontspec}
\IfFontExistsTF{Times New Roman}
  {\setmainfont{Times New Roman}}
  {\setmainfont{DejaVu Serif}}
\usepackage[ukrainian]{babel}
\usepackage{amsmath,amssymb,mathtools}
\usepackage{stmaryrd}
\usepackage{microtype}
\usepackage{enumitem}
\usepackage{booktabs}
\usepackage{array}
\usepackage[hidelinks,bookmarks=false]{hyperref}

\setlength{\parindent}{1.25em}
\setlength{\parskip}{0.25em}
\setlist[itemize]{leftmargin=1.6em}
\setlist[enumerate]{leftmargin=1.8em}

\newcommand{\ii}{\mathrm{i}}
\newcommand{\dd}{\mathop{}\!\mathrm{d}}
\newcommand{\jump}[1]{\llbracket #1\rrbracket}
\newcommand{\RePart}{\operatorname{Re}}

\title{\textbf{Незалежна перевірка співвідношень для розкриття\\
та енергетичного параметра зони передруйнування}}
\author{Технічна записка для обговорення}
\date{23 вересня 2026 р.}

\begin{document}

\maketitle

\section*{Мета записки}

Метою цієї записки є незалежна алгебраїчна перевірка двох співвідношень
поточного рукопису:
\begin{enumerate}
  \item формули для розкриття зони передруйнування $\delta_i$;
  \item формул для роботи сил опору відриву $W_i$ та пов'язаного з нею
  енергетичного параметра
  \[
    J_i=\frac{\dd W_i}{\dd d_i}.
  \]
\end{enumerate}

Перевірка виконується \emph{за умови прийняття наведеного в рукописі
розв'язку функціонального рівняння Вінера--Гопфа}. Тобто тут не
перевиводяться характеристичні визначники, ядра $G_i$ та сама процедура
факторизації; перевіряється лише подальше використання вже отриманого
$\Phi_i^-(p)$ при $p=0$ і $p=1$.

Такий спосіб дозволяє локалізувати можливу розбіжність і не змішувати її
з іншими питаннями аналітичного продовження або області збіжності
перетворення Мелліна.

\section{Вихідні співвідношення}

У поточному рукописі функціональне рівняння має вигляд
\begin{equation}
\label{eq:WH}
  \Phi_i^+(p)
  +\frac{\sigma_i}{p+1}
  -\frac{C Q_i d_i^\lambda}{p+1+\lambda}
  =
  -\tan(\pi p)\,G_i(p)\Phi_i^-(p),
\end{equation}
де $-1<\lambda<0$.

Використовується факторизація
\begin{equation}
\label{eq:Kdef}
  K^\pm(p)=
  \frac{\Gamma(1\mp p)}
       {\Gamma\!\left(\tfrac12\mp p\right)}
\end{equation}
та фактори $G_i^\pm(p)$, нормовані до одиниці на нескінченності.

Для правої трансформанти наведено розв'язок
\begin{equation}
\label{eq:PhiMinus}
\begin{aligned}
\Phi_i^-(p)
={}&K^-(p)G_i^-(p)
\Bigg[
\frac{\sigma_iK^+(-1)}
     {(p+1)G_i^+(-1)}
\\
&\qquad\qquad
-
\frac{C Q_i d_i^\lambda K^+(-1-\lambda)}
     {(1+\lambda)(p+1+\lambda)G_i^+(-1-\lambda)}
\Bigg].
\end{aligned}
\end{equation}

Для подальшого аналізу зручно ввести
\begin{equation}
\label{eq:AB}
  A_i=
  \frac{\sigma_iK^+(-1)}{G_i^+(-1)},
  \qquad
  B_i=
  \frac{C Q_i d_i^\lambda K^+(-1-\lambda)}
       {(1+\lambda)G_i^+(-1-\lambda)}.
\end{equation}
Тоді
\begin{equation}
\label{eq:PhiAB}
  \Phi_i^-(p)=
  K^-(p)G_i^-(p)
  \left[
    \frac{A_i}{p+1}
    -\frac{B_i}{p+1+\lambda}
  \right].
\end{equation}

\section{Умова регулярності на вершині зони}

При $p\to+\infty$ маємо
\[
  K^-(p)\sim p^{1/2},
  \qquad
  G_i^-(p)\to1,
\]
тому з \eqref{eq:PhiAB}
\[
  \Phi_i^-(p)
  \sim
  \frac{A_i-B_i}{\sqrt{p}}.
\]

З іншого боку, асимптотика біля вершини зони дає
\[
  \Phi_i^-(p)
  \sim
  -\frac{k_i}{\sqrt{2d_i p}}.
\]
Отже, умова відсутності кореневої сингулярності
\[
  k_i=0
\]
еквівалентна співвідношенню
\begin{equation}
\label{eq:AeqB}
  \boxed{A_i=B_i}.
\end{equation}

Саме з \eqref{eq:AeqB} випливає чинна формула для довжини зони
\begin{equation}
\label{eq:di}
  d_i=
  \left[
    \frac{C Q_iK^+(-1-\lambda)G_i^+(-1)}
         {\sigma_i(1+\lambda)K^+(-1)G_i^+(-1-\lambda)}
  \right]^{-1/\lambda}.
\end{equation}
Таким чином, подальша перевірка не ставить під сумнів
формулу \eqref{eq:di}.

Після підстановки \eqref{eq:AeqB} розв'язок
\eqref{eq:PhiAB} істотно спрощується:
\begin{align}
  \Phi_i^-(p)
  &=
  A_iK^-(p)G_i^-(p)
  \left[
    \frac{1}{p+1}
    -\frac{1}{p+1+\lambda}
  \right]
  \notag\\
  &=
  \boxed{
  A_iK^-(p)G_i^-(p)
  \frac{\lambda}
       {(p+1)(p+1+\lambda)}
  }.
\label{eq:PhiReduced}
\end{align}

Саме формула \eqref{eq:PhiReduced} є вихідною для двох незалежних
перевірок нижче.

\section{Перевірка формули для розкриття}

Розкриття зони в кутовій точці визначено як
\[
  \delta_i=\jump{u_\theta(0,\beta_i)}.
\]
За симетрії відносно бісектриси
$\delta_i=2u_\theta(0,\beta_i)$, а на зовнішньому кінці зони
$u_\theta(d_i,\beta_i)=0$. Тому безпосередньо з означення
$\Phi_i^-(p)$
\begin{align*}
  \Phi_i^-(0)
  &=
  \frac{E_i}{2(1-\nu_i^2)}
  \int_0^1
  \frac{\partial u_\theta}{\partial r}(\rho d_i,\beta_i)
  \,\dd\rho
  \\
  &=
  \frac{E_i}{2(1-\nu_i^2)d_i}
  \bigl[u_\theta(d_i,\beta_i)-u_\theta(0,\beta_i)\bigr]
  =
  -\frac{E_i\delta_i}{4(1-\nu_i^2)d_i}.
\end{align*}
Отже,
\begin{equation}
\label{eq:deltaPhi}
  \delta_i
  =
  -\frac{4(1-\nu_i^2)d_i}{E_i}\,
  \Phi_i^-(0).
\end{equation}

Покладемо $p=0$ у \eqref{eq:PhiReduced}. Тоді
\begin{equation}
\label{eq:Phi0}
  \Phi_i^-(0)
  =
  A_iK^-(0)G_i^-(0)
  \frac{\lambda}{1+\lambda}.
\end{equation}

Із \eqref{eq:Kdef}
\begin{equation}
\label{eq:K0}
  K^-(0)
  =
  \frac{\Gamma(1)}{\Gamma(1/2)}
  =
  \frac{1}{\sqrt{\pi}}.
\end{equation}

Для додатного парного ядра, факторизованого у прийнятому в рукописі
нормуванні,
\begin{equation}
\label{eq:Gsym}
  G_i^-(p)=\frac{1}{G_i^+(-p)}.
\end{equation}
Це співвідношення не є додатковим припущенням. Якщо позначити
\[
  \mathcal{L}_i(p)=
  \frac{1}{2\pi\ii}
  \int_{-\ii\infty}^{\ii\infty}
  \frac{\log G_i(z)}{z-p}\,\dd z,
\]
то для додатного парного ядра на осі факторизації можна вибрати
неперервну дійсну гілку $\log G_i(\ii t)$. Заміна $z\mapsto-z$
дає
\[
  \mathcal{L}_i(-p)=-\mathcal{L}_i(p),
\]
а отже, відповідно до означення факторів,
$G_i^+(-p)=1/G_i^-(p)$, що еквівалентно \eqref{eq:Gsym}.
Разом із граничним співвідношенням
$G_i(0)=G_i^+(0)/G_i^-(0)$ та вибором додатної гілки це дає
\begin{equation}
\label{eq:G0}
  G_i^-(0)=\frac{1}{\sqrt{G_i(0)}}.
\end{equation}
Значення факторів при $p=0$ тут і далі розуміються як відповідні
граничні значення з областей їх аналітичності.

Підставляючи \eqref{eq:AB}, \eqref{eq:K0} і \eqref{eq:G0}
у \eqref{eq:Phi0}, отримуємо
\begin{equation}
\label{eq:Phi0Final}
  \Phi_i^-(0)
  =
  \frac{\sigma_iK^+(-1)}
       {G_i^+(-1)}
  \frac{\lambda}
       {(1+\lambda)\sqrt{\pi G_i(0)}}.
\end{equation}

Отже, з \eqref{eq:deltaPhi}
\begin{equation}
\boxed{
  \delta_i=
  -\frac{4(1-\nu_i^2)\sigma_i\lambda K^+(-1)}
  {E_i\sqrt{\pi G_i(0)}(1+\lambda)\,G_i^+(-1)}
  d_i
  }.
\label{eq:deltaDerived}
\end{equation}

У поточному рукописі в знаменнику відповідної формули стоїть
\begin{equation}
\label{eq:deltaManuscript}
  G_i^+(-1-\lambda)
  \qquad\text{замість}\qquad
  G_i^+(-1).
\end{equation}

Таким чином, якщо виходити безпосередньо з
\eqref{eq:PhiMinus} та умови $k_i=0$, то аргумент
плюс-фактора у формулі розкриття має бути
\begin{equation}
  \boxed{-1}.
\end{equation}

Важливо, що $G_i^+(-1-\lambda)$ закономірно входить до
формули довжини \eqref{eq:di}: цей аргумент пов'язаний із полюсом
зовнішнього сингулярного поля $p=-1-\lambda$.
Натомість після використання умови $A_i=B_i$ значення
$\Phi_i^-(0)$ визначається коефіцієнтом
\[
  A_i=
  \frac{\sigma_iK^+(-1)}{G_i^+(-1)},
\]
тобто полюсом сталої напруги опору відриву при $p=-1$.

\section{Перевірка роботи $W_i$ та параметра $J_i$}

У поточному рукописі
\begin{equation}
\label{eq:Wdef}
  W_i=
  \sigma_i\int_0^{d_i}
  \jump{u_\theta(r,\beta_i)}\,\dd r.
\end{equation}

За симетрії $\jump{u_\theta}=2u_\theta$, а
$u_\theta(d_i,\beta_i)=0$. Тому інтегрування частинами дає
\[
  \int_0^{d_i}\jump{u_\theta}\,\dd r
  =
  -2\int_0^{d_i}
  r\frac{\partial u_\theta}{\partial r}\,\dd r.
\]
З іншого боку,
\[
  \Phi_i^-(1)
  =
  \frac{E_i}{2(1-\nu_i^2)d_i^2}
  \int_0^{d_i}
  r\frac{\partial u_\theta}{\partial r}\,\dd r.
\]
Отже,
\begin{equation}
\label{eq:WPhi}
  W_i=
  -\frac{4(1-\nu_i^2)\sigma_i}{E_i}
  d_i^2\Phi_i^-(1).
\end{equation}

Із \eqref{eq:PhiReduced} при $p=1$ маємо
\begin{equation}
\label{eq:Phi1}
  \Phi_i^-(1)
  =
  A_iK^-(1)G_i^-(1)
  \frac{\lambda}{2(2+\lambda)}.
\end{equation}

За \eqref{eq:Kdef}
\begin{equation}
\label{eq:Km1}
  K^-(1)
  =
  \frac{\Gamma(2)}{\Gamma(3/2)}
  =
  \frac{2}{\sqrt{\pi}},
  \qquad
  K^+(-1)
  =
  \frac{\Gamma(2)}{\Gamma(3/2)}
  =
  \frac{2}{\sqrt{\pi}}.
\end{equation}
Крім того, з \eqref{eq:Gsym}
\begin{equation}
\label{eq:G1}
  G_i^-(1)=\frac{1}{G_i^+(-1)}.
\end{equation}

Тому
\begin{equation}
\label{eq:Phi1Final}
  \Phi_i^-(1)
  =
  \frac{2\sigma_i\lambda}
  {\pi(2+\lambda)[G_i^+(-1)]^2}.
\end{equation}

Підстановка \eqref{eq:Phi1Final} у \eqref{eq:WPhi} дає
\begin{equation}
\boxed{
  W_i=
  -\frac{8\lambda(1-\nu_i^2)\sigma_i^2}
  {\pi(2+\lambda)[G_i^+(-1)]^2E_i}
  d_i^2
  }.
\label{eq:WDerived}
\end{equation}

Відповідно,
\begin{equation}
\boxed{
  J_i=
  \frac{\dd W_i}{\dd d_i}
  =
  -\frac{16\lambda(1-\nu_i^2)\sigma_i^2}
  {\pi(2+\lambda)[G_i^+(-1)]^2E_i}
  d_i
  }.
\label{eq:JDerived}
\end{equation}

У поточному рукописі наведено
\begin{align}
  W_i^{\mathrm{рук}}
  &=
  -\frac{32\lambda(1-\nu_i^2)\sigma_i^2}
  {9\pi(2+\lambda)[G_i^+(-1)]^2E_i}
  d_i^2,
\label{eq:WCurrent}
\\
  J_i^{\mathrm{рук}}
  &=
  -\frac{64\lambda(1-\nu_i^2)\sigma_i^2}
  {9\pi(2+\lambda)[G_i^+(-1)]^2E_i}
  d_i.
\label{eq:JCurrent}
\end{align}

Порівняння \eqref{eq:WDerived}--\eqref{eq:JDerived}
з \eqref{eq:WCurrent}--\eqref{eq:JCurrent} дає
\begin{equation}
\boxed{
  W_i=\frac94\,W_i^{\mathrm{рук}},
  \qquad
  J_i=\frac94\,J_i^{\mathrm{рук}}.
}
\label{eq:factor94}
\end{equation}

Для безрозмірного параметра
\[
  J_i'=
  \frac{E_iJ_i}
       {4(1-\nu_i^2)\sigma_i^2l}
\]
з \eqref{eq:JDerived} випливає
\begin{equation}
\boxed{
  J_i'
  =
  -\frac{4\lambda}
  {\pi(2+\lambda)[G_i^+(-1)]^2}
  \frac{d_i}{l}
  }.
\label{eq:JprimeDerived}
\end{equation}
У рукописі замість цього використано коефіцієнт
$-16/[9\pi]$.

\subsection*{Зауваження щодо числового коефіцієнта}

Коефіцієнт $32/(9\pi)$ можна точно відтворити, якщо замість
значення, що безпосередньо випливає з \eqref{eq:Kdef},
\[
  K^+(-1)=K^-(1)=\frac{2}{\sqrt{\pi}},
\]
підставити
\[
  \frac{4}{3\sqrt{\pi}}
  =
  \frac{\Gamma(2)}{\Gamma(5/2)}.
\]
Оскільки за означенням \eqref{eq:Kdef} у знаменнику при
$p=-1$ або $p=1$ має стояти $\Gamma(3/2)$, а не $\Gamma(5/2)$,
це спостереження є можливим поясненням появи множника $4/9$.
Водночас без первинних проміжних обчислень воно не доводить,
що саме так історично виник відповідний коефіцієнт.

\section{Числовий контроль для $\alpha=45^\circ$}

Для базового випадку
\[
  \frac{E_1}{E_2}=0.5,
  \qquad
  \nu_1=\nu_2=0.3,
  \qquad
  \alpha=45^\circ
\]
у поточному числовому розрахунку отримано
\[
  \lambda=-0.613197226573,
\]
\[
  G_1^+(-1-\lambda)=0.998663862615,
  \qquad
  G_1^+(-1)=0.996302154171.
\]
Ці значення плюс-факторів додатково відтворено прямим числовим
інтегруванням дійсної формули факторизації
\[
  \log G_1^+(p)
  =
  -\frac{p}{\pi}
  \int_0^\infty
  \frac{\log G_1(\ii t)}{t^2+p^2}\,\dd t,
  \qquad p<0.
\]

При $\sigma'=-0.5$ довжина зони
\[
  \frac{d_1}{l}=0.147434361703
\]
не змінюється, оскільки формула \eqref{eq:di} залишається тією самою.

За чинною формулою розкриття
\[
  \delta_1'^{\,\mathrm{рук}}
  =0.145442265256.
\]
Заміна лише аргументу плюс-фактора відповідно до
\eqref{eq:deltaDerived} дає
\[
  \boxed{
  \delta_1'=0.145787032378
  },
\]
тобто збільшення приблизно на $0.237\%$.

Для енергетичного параметра чинне значення
\[
  J_1'^{\,\mathrm{рук}}
  =0.037164616540
\]
за \eqref{eq:factor94} переходить у
\[
  \boxed{
  J_1'=0.083620387214
  }.
\]

Отже, виправлення формули розкриття дає порівняно невелику,
але залежну від кута зміну, тоді як виправлення енергетичного
коефіцієнта множить усю криву $J_i'$ на сталий множник $9/4$.

\section{Історична перевірка}

Було переглянуто чотири попередні публікації, пов'язані з тією самою
лінією аналітичних моделей.

\subsection*{Робота 2006 року}

У роботі A.~A.~Kaminsky, M.~V.~Dudik, L.~A.~Kipnis
(International Applied Mechanics, 42, 2006)~\cite{Kaminsky2006}
використано ту саму гамма-функціональну факторизацію
\[
  K^\pm(p)=
  \frac{\Gamma(1\mp p)}
       {\Gamma(\tfrac12\mp p)}.
\]
Робота присвячена зоні тангенціального розриву і не містить
спрощених формул для $\delta_i$, $W_i$ або $J_i$, які перевіряються
тут.

\subsection*{Робота 2007 року}

У роботі A.~A.~Kaminsky, M.~V.~Dudik, L.~A.~Kipnis
(International Applied Mechanics, 43, 2007)~\cite{Kaminsky2007}
розглянуто вже лінію
\emph{нормального} розриву переміщень. У точному розв'язку
$\Phi^-(p)$ внесок сталої напруги опору відриву містить
\[
  \frac{\sigma K^+(-1)}
       {(p+1)G^+(-1)},
\]
тоді як полюси зовнішнього осцилюючого поля зберігають власні
аргументи факторів. Така структура узгоджується з розділенням
аргументів $-1$ і $-1-\lambda$ у наведеному вище незалежному
виведенні.

У тій самій роботі енергія визначається через момент
$\Phi^-(1)$, що безпосередньо підтверджує використання
$p=1$ для енергетичної характеристики. Окремої формули розкриття
через $\Phi^-(0)$ у цій роботі не вводиться, тому робота 2007 року
не використовується тут як підтвердження саме кроку $p=0$.

Слід, однак, враховувати, що в роботі 2007 року використано інше
визначення енергетичної величини з множником $1/2$:
\[
  W=\frac{\sigma}{2}\int \jump{u_\theta}\,\dd r.
\]
Тому її числовий коефіцієнт не слід безпосередньо переносити
до поточного визначення \eqref{eq:Wdef}.

\subsection*{Робота 2023 року}

У роботі A.~Kaminsky, M.~Dudyk, Y.~Reshitnyk, Y.~Chornoivan
(International Journal of Solids and Structures, 267, 2023)
~\cite{Kaminsky2023} використано ту саму факторизацію
\[
  K^\pm(p)=
  \frac{\Gamma(1\mp p)}
       {\Gamma(\tfrac12\mp p)}.
\]
У ній енергетична характеристика явно виражається через
трансформанти $\Phi_m^-(1)$, а стрибки переміщень у вершині
тріщини --- через $\Phi_m^-(0)$. Отже, ця робота незалежно
підтверджує сам вибір моментів $p=1$ та $p=0$ для відповідних
фізичних величин, хоча її векторна задача не визначає числових
коефіцієнтів поточної скалярної моделі.

\subsection*{Робота 2024 року}

У роботі A.~O.~Kaminsky, M.~V.~Dudyk, T.~V.~Polishchuk
(International Applied Mechanics, 60, 2024)~\cite{Kaminsky2024}
уже присутні обидва структурні елементи, що перевіряються:
\begin{itemize}
  \item у формулі розкриття використано
  $G_i^+(-1-\lambda)$;
  \item у формулах для $W_i$ та $J_i$ використано коефіцієнти
  $32/(9\pi)$ та $64/(9\pi)$ відповідно.
\end{itemize}

Отже, ці дві розбіжності не виникли під час підготовки поточного
рукопису: відповідний аргумент плюс-фактора та коефіцієнти вже
були в опублікованих формулах 2024 року. Окремо зауважимо, що
в друкованих формулах 2024 року для $W_i$ та $J_i$ міститься
лише один множник $\sigma_i$, тоді як у поточному рукописі
розмірнісно необхідний множник $\sigma_i^2$ уже відновлено.
Це окреме виправлення не є предметом теперішньої перевірки.
Водночас у роботах 2006--2007 років не знайдено відповідних
спрощених співвідношень, які б вимагали саме таких аргументів
або коефіцієнтів.

\section{Питання для авторського рішення}

На підставі наведеного виведення пропонується перевірити такі два
пункти.

\begin{enumerate}
  \item Чи слід у формулі розкриття замінити
  \[
    \boxed{
    G_i^+(-1-\lambda)\;\longrightarrow\;G_i^+(-1)
    }?
  \]

  \item Чи слід замінити
  \[
    \boxed{
    \frac{32}{9\pi}\;\longrightarrow\;\frac{8}{\pi}
    }
    \qquad\text{у }W_i,
  \]
  та відповідно
  \[
    \boxed{
    \frac{64}{9\pi}\;\longrightarrow\;\frac{16}{\pi}
    }
    \qquad\text{у }J_i?
  \]
\end{enumerate}

Якщо авторське виведення приводить до іншого результату, корисно
уточнити додаткове співвідношення, інше нормування факторів або іншу
інтерпретацію $\Phi_i^-(p)$, яка змінює перехід від
\eqref{eq:PhiMinus} до значень при $p=0$ та $p=1$.

До остаточного авторського рішення доцільно не змінювати
формулу довжини зони та не перераховувати рисунки. Після
підтвердження:
\begin{itemize}
  \item крива $d_i/l$ залишається незмінною;
  \item криві $\delta_i'$ слід перерахувати з
  $G_i^+(-1)$;
  \item криві $J_i'$ слід помножити на $9/4$;
  \item критичні навантаження, визначені з енергетичного критерію
  $J_i=J_{ic}$, потребуватимуть окремого перерахунку.
\end{itemize}

\begin{thebibliography}{9}

\bibitem{Kaminsky2006}
A.~A.~Kaminsky, M.~V.~Dudik, and L.~A.~Kipnis,
``On the Direction of Development of a Thin Fracture Process Zone
at the Tip of an Interfacial Crack Between Dissimilar Media,''
\emph{International Applied Mechanics}, vol.~42, no.~2,
pp.~136--144, 2006.
\href{https://doi.org/10.1007/s10778-006-0068-1}
{doi:10.1007/s10778-006-0068-1}.

\bibitem{Kaminsky2007}
A.~A.~Kaminsky, M.~V.~Dudik, and L.~A.~Kipnis,
``Initial Kinking of an Interface Crack Between Two Elastic Media,''
\emph{International Applied Mechanics}, vol.~43, no.~10,
pp.~1090--1099, 2007.
\href{https://doi.org/10.1007/s10778-007-0109-4}
{doi:10.1007/s10778-007-0109-4}.

\bibitem{Kaminsky2023}
A.~Kaminsky, M.~Dudyk, Y.~Reshitnyk, and Y.~Chornoivan,
``An analytical method of modeling the process zone near the tip of
an interface crack due to its kinking from the interface of
quasi-elastic materials,''
\emph{International Journal of Solids and Structures},
vol.~267, art.~112117, 2023.
\href{https://doi.org/10.1016/j.ijsolstr.2023.112117}
{doi:10.1016/j.ijsolstr.2023.112117}.

\bibitem{Kaminsky2024}
A.~O.~Kaminsky, M.~V.~Dudyk, and T.~V.~Polishchuk,
``Crack Initiation Model in Elastic Piecewise Homogeneous Body
with Broken Interface,''
\emph{International Applied Mechanics}, vol.~60, no.~6,
pp.~677--688, 2024.
\href{https://doi.org/10.1007/s10778-025-01319-8}
{doi:10.1007/s10778-025-01319-8}.

\end{thebibliography}

\end{document}